\documentclass[fontset=none,11pt,a4paper]{ctexart}
\setCJKmainfont{AR PL SungtiL GB}[BoldFont=AR PL UMing CN,AutoFakeBold=2.5]
\setCJKsansfont{AR PL UMing CN}
\setCJKmonofont{AR PL UMing CN}
\usepackage{geometry,booktabs,multirow,hyperref,enumitem,xcolor,longtable}
\geometry{margin=2.2cm}
\title{多模态Agent Post-Training项目\\完整分析报告}
\author{}
\date{\today}

\begin{document}
\maketitle

\begin{abstract}
本项目以Qwen2.5-VL-3B/7B为基座，训练VLM在发票核验、图表计算、图文一致性等场景下严格遵循vision\_parse $\to$ retrieve\_docs/python\_exec $\to$ final\_answer的多步工具调用流程。本文从代码架构、问题演化、优化历程、失败分析、当前瓶颈、后续方向六个维度进行系统性梳理。核心发现：(1) 单阶段SFT在同时优化格式纪律、工具选择、答案质量、拒答边界时存在根本性多目标冲突；(2) 项目中唯一有效的优化全部来自多模态层面（真实VLM替代mock、3B升7B），Agent层面的增量修补（DPO、增强数据、微调）全部失败；(3) 7B在文档理解上碾压3B，但对训练数据偏移的敏感度呈数量级增加；(4) 当前最优为V4(3B)和V7B-clean(7B)，后续RL优化的真正空间在7B的工具纪律修复。
\end{abstract}

\tableofcontents

\section{项目定义}

\subsection{任务描述}

给定一张图片和一个自然语言问题，VLM需按以下流程完成任务：
\begin{enumerate}[nosep]
    \item 判断任务类型（发票核验/图表计算/图文一致性/品牌识别/简单描述）
    \item 如需视觉信息，调用 vision\_parse(mode) 进行OCR/图表读数/图像描述
    \item 如需规则依据，调用 retrieve\_docs(query) 检索规则文档
    \item 如需数学计算，调用 python\_exec(code) 完成计算
    \item 输出包含依据的 final\_answer 或合理拒答
\end{enumerate}

核心约束：每轮对话模型只能输出一个 \texttt{<tool\_call>} 或一个 \texttt{<final\_answer>}，不得同时输出多个动作或裸文本。这一格式约束是整个项目的阿喀琉斯之踵。

\subsection{代码架构}

\begin{longtable}{lp{9cm}}
\toprule
文件 & 功能 \\
\midrule
agent/runtime.py & 多轮Agent执行引擎---管理消息循环、enforce拦截、工具调用 \\
agent/tools.py & 三个工具实现---vision\_parse(真实VLM/mock)、retrieve\_docs(规则检索)、python\_exec(安全eval) \\
agent/vlm.py & Qwen-VL模型封装---加载、LoRA注入、mock回退 \\
agent/parser.py & 输出解析---正则提取tool\_call或final\_answer，容错缺闭合标签 \\
agent/prompts.py & System prompt---定义工具调用格式规范 \\
\midrule
train/train\_sft\_lora.py & SFT LoRA训练---expand\_assistant\_turns + prefix/full collator \\
train/train\_dpo\_lora.py & DPO偏好训练---last-turn log-prob优化（三次崩溃，已弃用） \\
train/train\_rl\_custom.py & 自研RL训练---REINFORCE+KL，无TRL依赖（当前迭代中） \\
\midrule
eval/run\_vlm\_eval.py & 评测主脚本---ToolHit, ToolOrder, Format, Finished, Keyword, Refusal \\
eval/run\_seed\_eval.py & 早期简化评测（不依赖真实VLM） \\
\midrule
scripts/build\_agent\_train\_real\_v3.py & 合并SFT+外部数据集，控制上采样和recovery参数 \\
scripts/rebuild\_seed\_with\_real\_vlm.py & 用真实VLM重建训练数据中的vision\_parse结果 \\
scripts/real\_vision\_parse.py & 共享真实VLM推理模块 \\
scripts/convert\_chartqa/docvqa/cord\_to\_sft.py & 外部数据集转换为Agent SFT格式 \\
scripts/build\_augmented\_v5\_data.py & 数据增强---多样问法、推理链、纠错轨迹 \\
scripts/build\_hard\_cases\_from\_eval.py & 从评测错误生成SFT修正样本和DPO偏好对 \\
scripts/build\_tool\_decision\_samples.py & 工具决策样本---直接回答+refusal+tool \\
scripts/build\_refusal\_samples.py & 用7B真实VLM生成防幻觉样本 \\
scripts/tag\_vlm\_errors.py & 评测结果自动错误标注 \\
\bottomrule
\end{longtable}

\subsection{评测体系}

50条评测样本，分5类：document\_validation(15)、image\_text\_consistency(12)、chart\_calculation(10)、uncertainty\_refusal(8)、direct\_vqa(5)。

6项指标：
\begin{enumerate}[nosep]
    \item Tool Hit：gold\_tools $\subseteq$ pred\_tools
    \item Tool Order Hit：gold\_tools顺序一致
    \item Format Valid：无parse\_error
    \item Finished：正常结束（未超max\_steps）
    \item Answer Keyword Hit：gold\_keywords全部出现在答案中
    \item Refusal Correct：should\_refuse与is\_refusal一致
\end{enumerate}

\section{分阶段历程}

\subsection{阶段0：基础建设（V2之前）}

\textbf{问题}：模型不会按流程调工具。

\textbf{做法}：手工构造50条seed数据，每条包含完整的多轮工具调用轨迹。通过build\_first\_tool\_sft.py、build\_second\_tool\_sft.py、build\_sft\_train.py三个脚本拆分为first\_tool、full\_trace、second\_tool三类训练样本。

\subsection{阶段1：Mock数据时代（V2）}

\textbf{问题}：工具命中率只有50\%，模型经常跳过第二工具。

\textbf{做法}：在runtime.py实现enforce\_required\_tools---关键词匹配判断是否需要第二工具，漏了就拦截final\_answer并注入纠正提示。加入ChartQA外部数据（446条）。

\textbf{关键代码}（runtime.py）：
\begin{verbatim}
required_tool = self._required_tool_after_vision(question, called_tools)
if self.enforce_required_tools and required_tool:
    messages.append({"role": "user",
        "content": self._missing_required_tool_prompt(required_tool)})
    continue
\end{verbatim}

\textbf{效果}：工具命中率从50\%升至100\%。但这是runtime兜底，不是模型学会的。

\subsection{阶段2：V3过度工具调用}

\textbf{问题}：模型在不需要工具的场景也强行调工具。根因是recovery变体教模型"宁可多调不能漏调"，ChartQA/CORD上采样5x过度放大第二工具信号。

\textbf{做法}：recovery\_repeat 2$\to$1，python\_exec/retrieve\_docs repeat 5$\to$2，recovery提示词对齐runtime。

\subsection{阶段3：真实VLM数据（V4）---全项目最大单次提升}

\textbf{问题}：训练数据中vision\_parse结果是人工mock的，外部数据集甚至用占位文本。模型从未见过真实VLM输出分布。

\textbf{做法}：
\begin{itemize}[nosep]
    \item 新建rebuild\_seed\_with\_real\_vlm.py：加载真实Qwen2.5-VL-3B，对46张种子图片逐张做vision\_parse，替换mock结果
    \item 新建real\_vision\_parse.py：共享真实VLM推理模块
    \item 修改convert\_chartqa/docvqa/cord\_to\_sft.py：支持--use\_real\_vlm标志
    \item 797张图片全部完成真实VLM处理
    \item 外部图片Resize到max 1024px（DocVQA最大1.8MB、CORD最大7.4MB导致OOM）
\end{itemize}

\textbf{效果}：V4成为3B最优模型。这是全项目最大单次提升。

\subsection{阶段4：数据增强（V5/V6/V8）---全部失败或持平}

\textbf{做法}：
\begin{itemize}[nosep]
    \item V5：build\_augmented\_v5\_data.py生成818条增强数据（632条多样问法+80条推理链+42条纠错轨迹+6条refual+8条直接回答）
    \item V6：修复refual模板（从假模板换为真实VLM输出）
    \item V8：build\_tool\_decision\_samples.py生成53条工具决策样本
\end{itemize}

\textbf{结果}：V5持平(82\%)、V6退步(70\%)、V8大幅退步(50\%)。

\textbf{根因}：单阶段SFT需要同时优化四个互相矛盾的目标---格式纪律、工具选择、答案质量、拒答边界。这些在梯度空间中是竞争关系。任何针对单一目标的修补都会打破平衡。

\subsection{阶段5：DPO（三次尝试）---全部崩溃}

\textbf{现象}：
\begin{itemize}[nosep]
    \item DPO单轮（64对）：tool hit 100\%$\to$60\%
    \item DPO多轮（13对）：tool hit$\to$14\%，输出空字符串
    \item SFT微调（13条）：tool hit$\to$22\%
\end{itemize}

\textbf{根因}：train\_dpo\_lora.py的compute\_log\_prob只优化最后一条assistant turn。偏好数据是单轮的user$\to$final\_answer。模型可以在降低DPO loss的同时完全抛弃多轮格式约束。

\subsection{阶段6：7B迁移（V7B-clean）---文档理解质的飞跃}

\textbf{问题}：3B的视觉理解能力是瓶颈。

\textbf{做法}：用与V4完全相同的1133条数据，换Qwen2.5-VL-7B-Instruct。经过rank搜索：r=16崩溃(14\% tool)、r=32成功(92\%)、r=64崩溃(16\%)。

\textbf{关键发现}：7B对训练数据分布偏移的敏感度比3B高一个数量级。加9条refusal样本（0.8\%的数据改动）就能让工具命中从92\%崩到16\%。

\textbf{效果}：document\_validation 10/15$\to$15/15，image\_text\_consistency 9/12$\to$12/12。

\subsection{阶段7：RL/GRPO（当前迭代）}

经历了四轮迭代：
\begin{enumerate}[nosep]
    \item 用TRL的GRPOTrainer---失败（Qwen2.5-VL RoPE索引兼容性bug）
    \item 手写GRPO，ratio分母=$\pi_{ref}$（SFT）---3B跑11轮reward不动，模型被困在SFT局部最优
    \item 修正ratio分母=$\pi_{old}$，用跨轮缓存---loss爆炸(1.58$\to$4.51)，根因是缓存按位置匹配但文本每轮不同，ratio成纯噪声
    \item 放弃跨轮ratio，改用纯REINFORCE+KL（ratio=1），kl\_coef$\uparrow$0.2, lr$\downarrow$1e-6---当前版本
\end{enumerate}

\section{核心问题总结}

\begin{longtable}{p{0.5cm}p{3.5cm}p{3.5cm}p{3cm}p{1.5cm}}
\toprule
\# & 问题 & 根因 & 解决 & 效果 \\
\midrule
1 & Mock数据 vs 真实VLM不匹配 & 训练用人工OCR结果 & rebuild\_seed\_with\_real\_vlm.py重建797张图 & 全项目最大提升 \\
2 & 大图OOM & Qwen-VL动态分辨率 & Resize到max 1024px & 训练可行 \\
3 & SFT多目标冲突 & 格式+工具+答案+拒答抢一个loss & 未解决 & 所有增量补丁失败 \\
4 & DPO格式崩溃 & DPO只优化last-turn & 放弃DPO & 三次崩溃 \\
5 & 7B数据极度敏感 & 大模型捷径学习 & 只用纯V4数据，不添加额外样本 & r=32成功 \\
6 & GRPO跨轮ratio噪声 & 缓存按位置匹配，文本每轮不同 & 改为ratio=1（纯REINFORCE） & 待验证 \\
7 & 7B双模型GPU OOM & 两个7B$\approx$30GB & --ref\_on\_cpu flag & 待验证 \\
8 & 3B reward到天花板 & V4已100\%tool & GRPO对3B无效 & 放弃3B，聚焦7B \\
\bottomrule
\end{longtable}

\section{优化上限分析}

\subsection{3B天花板}

V4(3B)的工具纪律已达100\%，文档理解受限于3B OCR质量。群体rollout reward同质化，GRPO无优化信号。3B已到极限。

\subsection{7B天花板}

V7B-clean(7B)文档理解满分(15/15+12/12)，但工具纪律92\%。品牌识别存在真实的漏调（refusal\_002的"Vans"幻觉），reward有区分度。

\textbf{7B是RL的真正目标。}

\subsection{RL理论上限}

\begin{enumerate}[nosep]
    \item REINFORCE固有方差问题：40样本$\times$4rollout的梯度噪声大，需1000+ rollout才能稳定收敛
    \item KL锚定SFT限制了探索范围---tradeoff安全性 vs 优化空间
    \item reward函数设计决定了优化方向：工具奖励(1.0)$>$答案奖励(0.5)$>$拒绝奖励(0.5)$>$惩罚(-2.0$\sim$-0.3)
\end{enumerate}

\section{当前最优模型}

\begin{table}[htbp]
\centering
\begin{tabular}{lcc}
\toprule
指标 & V4 (3B) & V7B-clean (7B) \\
\midrule
Tool Hit & \textbf{50/50} & 46/50 (无enf 49/50) \\
Keyword & 38/50 & \textbf{43/50} \\
文档核验 & 10/15 & \textbf{15/15} \\
图文一致性 & 9/12 & \textbf{12/12} \\
品牌识别 & 8/8 tool & 4/8 tool \\
数据敏感性 & \textbf{低} & 极高 \\
RL优化空间 & 无（天花板） & \textbf{有（92\%$\to$100\% tool）} \\
\bottomrule
\end{tabular}
\caption{两个最优模型的能力互补}
\end{table}

\section{RL训练命令}

\begin{verbatim}
# 3B（验证用，两个模型都在GPU）
python train/train_rl_custom.py \
    --model_name .../Qwen2.5-VL-3B-Instruct \
    --adapter_name .../lora_agent_real_v4 \
    --output_dir .../lora_grpo_v1 --bf16

# 7B（主跑，ref放CPU避免OOM）
python train/train_rl_custom.py \
    --model_name .../Qwen2.5-VL-7B-Instruct \
    --adapter_name .../lora_agent_v7b_clean \
    --output_dir .../lora_grpo_v2_7b \
    --ref_on_cpu --lr 1e-6 --kl_coef 0.2
\end{verbatim}

\section{后续方向}

\begin{enumerate}[nosep]
    \item 7B + RL优化：修复品牌识别(92\%$\to$100\% tool)，消除"Vans"幻觉
    \item reward函数增强：对parse\_error加重罚分(-2$\to$-5)，强化格式保护
    \item 增大rollout数量：K=4$\to$8，增加采样多样性
    \item 跨模型蒸馏：V4的工具纪律$\to$V7B，融合两者优势
    \item 模型规模工具学习动力学研究：为什么7B文档理解强但工具纪律弱？为什么7B对数据偏移更敏感？
\end{enumerate}

\end{document}
